\documentclass[11pt,letterpaper]{article}
\usepackage{jahnel}
\usepackage{framed}
\DeclareMathOperator{\diam}{diam}
\DeclareMathOperator{\dist}{dist}
\newtheorem{claim}{Claim}
\crefname{claim}{Claim}{Claims}
\Crefname{claim}{Claim}{Claims}
\newenvironment{boxedclaim}[2]{\renewcommand{\theclaim}{#1}\begin{framed}\begin{claim}[#2]}{\end{claim}\end{framed}}
\newcommand{\cS}{\mathcal S}
\renewcommand{\cE}{\mathcal E}
\newcommand{\cR}{\mathcal R}

\runningtitle{A two-scale scheme for the Cox Boolean model}
\hypersetup{pdftitle={A two-scale scheme for the Cox Boolean model on Poisson lines},pdfauthor={Christian Hirsch, Benedikt Jahnel}}

\title{\textbf{A two-scale scheme for the Cox Boolean model on Poisson lines}\\[0.6em] \large Working note}

\author{%
  Christian Hirsch\thanks{Institut f\"ur Mathematische Stochastik, Georg-August-Universit\"at G\"ottingen, Goldschmidtstra\ss e~7, 37077 G\"ottingen, Germany. \texttt{christian.hirsch@uni-goettingen.de}. \orcid{0000-0003-4136-3740}.}
  \and
  Benedikt Jahnel\thanks{Institut f\"ur Mathematische Stochastik, Technische Universit\"at Braunschweig, Universit\"atsplatz~2, 38106 Braunschweig, Germany, and Weierstrass Institute for Applied Analysis and Stochastics, Mohrenstra\ss e~39, 10117 Berlin, Germany. \texttt{benedikt.jahnel@tu-braunschweig.de}. \orcid{0000-0002-4212-0065}.}%
}

\date{\today}

\begin{document}

\maketitle
\thispagestyle{plain}

\begin{abstract}
We adapt Gou\'er\'e's two-scale argument for the subcritical phase of the Poisson Boolean model to the Boolean model driven by a Cox process on a Poisson line process. The argument is run directly on occupied crossings of annuli, so that neither thick clusters, nor the skeleton bound, nor the first moment over cluster chains of our earlier notes is needed. Two things change relative to the Poisson Boolean model. The inner and outer regions of the original argument are replaced by two bounded regions of the same size, because every line through a bounded region reaches its complement; and the dependence between the two regions, which in the original argument is carried by rare large balls, is carried here by the lines meeting both regions, a constant fraction of all lines through either, and is handled by conditioning on them and treating them as fixed lines. The induction hypothesis therefore carries fixed lines with a weight, and we show that the recursion then becomes a deterministic inequality for the pair (target function, weight) in which randomness enters only through the law of the shared lines. Defining the weight recursively makes the crossing bound unconditional and reduces the existence of a subcritical phase to a single statement, that the recursively defined target function tends to zero. We prove the inclusion lemma, the decoupling, the bounds on shared lines, the recursion theorem, and the closure of one recursion step under a group-weight hypothesis on the shared lines, and we identify what the hypothesis must look like for it to propagate. We also correct a claim of our multiscale note: a polynomial or faster target does not force many sub-boxes at polynomially separated scales; two regions at a fixed ratio suffice, because the quadratic term of the recursion is far below the target.
\end{abstract}

% =================================================================
\section{Setting}
\label{sec:setting}

Let $\cL$ be a Poisson line process in~$\R^2$ with intensity measure $\nu(\dd\theta\,\dd p)=(\mu/\pi)\,\dd\theta\,\dd p$, so that the lines meeting a convex set of perimeter~$P$ form a Poisson variable with mean $\mu P/\pi$; the lines meeting a disk of radius~$r$ have measure $2\mu r$. For a finite set~$G$ of lines (\emph{fixed lines}), let $X$ be, conditionally on~$\cL$, a Poisson process of intensity~$\lambda$ per unit length on every line of $\cL\cup G$. For a Borel set~$A$ let
\begin{equation*}
  \Xi(A)\coloneqq\bigcup_{c\in X\cap A}B(c,1)
\end{equation*}
be the union of the unit disks centred at the points of~$X$ in~$A$, and $\Xi\coloneqq\Xi(\R^2)$. We write $\P_G$ for the law of the model with fixed lines~$G$ and $\P=\P_\varnothing$. The radius~$1$ is no restriction by scaling.

\paragraph{Crossing events.} Following~\cite{Gou08}, for $x\in\R^2$ and $\alpha\ge1$ let $\cE(x,\alpha)$ be the event that the connected component of
\begin{equation*}
  \Xi\bigl(B(x,10\alpha)\bigr)\cup B(x,\alpha)
\end{equation*}
containing~$x$ is not contained in $B(x,8\alpha)$. The event depends only on the lines of $\cL\cup G$ meeting $B(x,10\alpha)$ and on the points of~$X$ in $B(x,10\alpha)$. Put
\begin{equation}
\label{eq:pi}
  \pi_\alpha(x;G)\coloneqq\P_G\bigl(\cE(x,\alpha)\bigr)\,,\qquad\pi_\alpha\coloneqq\pi_\alpha(0;\varnothing)\,.
\end{equation}
Without fixed lines the model is stationary and $\pi_\alpha(x;\varnothing)=\pi_\alpha$.

\begin{lemma}[Percolation forces crossings]
\label{lem:perc}
Let $\alpha\ge1$. If the component of~$\Xi$ containing a point of $B(0,1)$ is unbounded, then $\cE(0,\alpha)$ occurs. Consequently, if $\pi_\alpha\to0$ as $\alpha\to\infty$, then $\Xi$ has no unbounded component almost surely.
\end{lemma}

\begin{proof}
Let $D_1,D_2,\dots$ be a chain of pairwise consecutively intersecting disks of~$\Xi$ with $D_1\cap B(0,1)\neq\varnothing$ that leaves $B(0,9\alpha)$; such a chain exists if the component is unbounded, since the component is a union of disks and is path connected through consecutive intersections. Let $D_n$ be the first disk of the chain not contained in $B(0,8\alpha)$. The centres of $D_1,\dots,D_{n-1}$ lie in $B(0,8\alpha)$, and the centre of~$D_n$ is within distance~$1$ of a point of $B(0,8\alpha)$, since $D_n$ meets $D_{n-1}$ or, if $n=1$, meets $B(0,1)$; so all these centres lie in $B(0,10\alpha)$ for $\alpha\ge1$. The set $B(0,\alpha)\cup D_1\cup\dots\cup D_n$ is connected, contains~$0$ and is not contained in $B(0,8\alpha)$, which is $\cE(0,\alpha)$. For the consequence, if an unbounded component existed with positive probability, by stationarity it would meet $B(z,1)$ for some $z\in\Z^2$ with positive probability, hence $B(0,1)$ with positive probability, and then $\pi_\alpha\ge\P(\text{unbounded component meets }B(0,1))>0$ for every~$\alpha$.
\end{proof}

% =================================================================
\section{The two-scale inclusion}
\label{sec:inclusion}

Fix finite sets $K\subseteq S_{10}$ and $L\subseteq S_{80}$ of points on the circles of radii $10$ and $80$ around the origin such that $S_{10}\subseteq K+B(0,1)$ and $S_{80}\subseteq L+B(0,1)$; one can take $\lvert K\rvert\le70$ and $\lvert L\rvert\le510$. For $x\in\R^2$, $\alpha\ge1$, $k\in K$ and $l\in L$ write
\begin{equation*}
  x_k\coloneqq x+\alpha k,\qquad x_l\coloneqq x+\alpha l,\qquad R_k\coloneqq B(x_k,10\alpha),\qquad R_l\coloneqq B(x_l,10\alpha)\,.
\end{equation*}
Then $R_k\subseteq B(x,20\alpha)$, $R_l\subseteq B(x,70\alpha)^c$, and $\lvert x_k-x_l\rvert\ge70\alpha$.

\begin{lemma}[Inclusion, after {\cite[Lemma~3.3]{Gou08}}]
\label{lem:inclusion}
For every finite set~$G$ of lines, every $x\in\R^2$ and $\alpha\ge1$,
\begin{equation}
\label{eq:inclusion}
  \cE(x,10\alpha)\subseteq\bigcup_{k\in K}\bigcup_{l\in L}\bigl(\cE(x_k,\alpha)\cap\cE(x_l,\alpha)\bigr)\,.
\end{equation}
\end{lemma}

In~\cite{Gou08} the inclusion holds only outside the event that a ball of radius at least~$\alpha$ is present, and that event supplies the error term of the recursion. Here all disks have radius $1\le\alpha$ and the inclusion holds outright.

\begin{proof}
Assume $\cE(x,10\alpha)$ and take $x=0$ without loss of generality. The component of $\Xi(B(0,100\alpha))\cup B(0,10\alpha)$ containing~$0$ is not contained in $B(0,80\alpha)$; since $X\cap B(0,100\alpha)$ is finite almost surely, there is a chain $D_1,\dots,D_n$ of unit disks with centres in $B(0,100\alpha)$, consecutive disks intersecting, with $D_1\cap B(0,10\alpha)\neq\varnothing$ and $D_n\not\subseteq B(0,80\alpha)$. The union $U\coloneqq D_1\cup\dots\cup D_n$ is connected and contains a point of $B(0,10\alpha)$ and a point outside $B(0,80\alpha)$, so it contains points $y\in S_{10\alpha}$ and $y'\in S_{80\alpha}$; say $y\in D_i$ and $y'\in D_m$. Choose $k\in K$ with $\lvert y-\alpha k\rvert<\alpha$ and $l\in L$ with $\lvert y'-\alpha l\rvert<\alpha$, so that $D_i$ meets $B(x_k,\alpha)$ and $D_m$ meets $B(x_l,\alpha)$.

\emph{The event at~$k$.} The union $D_i\cup\dots\cup D_n$ is connected, meets $B(x_k,\alpha)$ and contains a point outside $B(0,80\alpha)\supseteq B(x_k,8\alpha)$. Let $D_{i'}$, $i'\ge i$, be the first of these disks not contained in $B(x_k,8\alpha)$. The centres of $D_i,\dots,D_{i'-1}$ lie in $B(x_k,8\alpha)$, and the centre of $D_{i'}$ is within distance~$1$ of a point of $B(x_k,8\alpha)$, because $D_{i'}$ meets $D_{i'-1}$ if $i'>i$ and meets $B(x_k,\alpha)$ if $i'=i$; so all these centres lie in $R_k=B(x_k,10\alpha)$. Hence $B(x_k,\alpha)\cup D_i\cup\dots\cup D_{i'}$ is a connected subset of $\Xi(R_k)\cup B(x_k,\alpha)$ containing~$x_k$ and not contained in $B(x_k,8\alpha)$: this is $\cE(x_k,\alpha)$.

\emph{The event at~$l$.} The union $D_1\cup\dots\cup D_m$ is connected, meets $B(x_l,\alpha)$ and contains a point of $B(0,10\alpha)$, which lies outside $B(x_l,8\alpha)$ since $\lvert x_l\rvert=80\alpha$. Running the chain backwards from~$D_m$, the same argument gives $\cE(x_l,\alpha)$.
\end{proof}

% =================================================================
\section{Decoupling through the shared lines}
\label{sec:decoupling}

For disjoint bounded Borel sets $A_1,A_2$ let $\cS=\cS(A_1,A_2)$ be the set of lines of~$\cL$ meeting both, and $\mathcal P_i$ the set of lines of~$\cL$ meeting $A_i$ but not the other set (the \emph{private} lines of~$A_i$). For a finite set~$H$ of lines write $H\cap A\coloneqq\{\ell\in H\colon\ell\cap A\neq\varnothing\}$.

\begin{lemma}[Two levels of independence]
\label{lem:decoupling}
Let $A_1,A_2$ be disjoint bounded Borel sets, $G$ a finite set of lines, and let $F_i$ be an event determined by the lines of $\cL\cup G$ meeting~$A_i$ and the points of~$X$ in~$A_i$, $i=1,2$. Then
\begin{equation}
\label{eq:decoupling}
  \P_G(F_1\cap F_2)=\E\Bigl[\P_{G\cup\cS}(F_1)\,\P_{G\cup\cS}(F_2)\Bigr]\,,
\end{equation}
where inside the expectation $\cS$ is treated as a deterministic set of fixed lines and $\P_{G\cup\cS}(F_i)$ is the probability of~$F_i$ in the model whose random lines are the Poisson process of lines meeting $A_i$ but not $A_{3-i}$ and whose fixed lines are $(G\cup\cS)\cap A_i$.
\end{lemma}

\begin{proof}
The sets of lines meeting $A_1$ only, meeting $A_2$ only, and meeting both are disjoint, so $\mathcal P_1,\mathcal P_2,\cS$ are independent Poisson processes. Conditionally on all lines, the points of~$X$ in $A_1$ and in $A_2$ are independent Poisson processes on the lines, since $A_1\cap A_2=\varnothing$. Hence, conditionally on~$\cS$, the pairs $(\mathcal P_1,X\cap A_1)$ and $(\mathcal P_2,X\cap A_2)$ are independent, and $F_i$ is determined by $(\mathcal P_i,X\cap A_i)$ together with the fixed lines $(G\cup\cS)\cap A_i$. Taking the conditional expectation given~$\cS$ gives~\eqref{eq:decoupling}.
\end{proof}

Applied to the inclusion~\eqref{eq:inclusion} with $A_1=R_k$, $A_2=R_l$,
\begin{equation}
\label{eq:step}
  \pi_{10\alpha}(x;G)\le\sum_{k\in K}\sum_{l\in L}\E\Bigl[\pi_\alpha\bigl(x_k;G\cup\cS_{kl}\bigr)\,\pi_\alpha\bigl(x_l;G\cup\cS_{kl}\bigr)\Bigr],\qquad\cS_{kl}\coloneqq\cS(R_k,R_l)\,,
\end{equation}
where only the lines of $G\cup\cS_{kl}$ meeting the respective region matter in each factor. This is the two-scale step. Its content is that the two sub-crossings are independent given the shared lines, and that the shared lines enter as fixed lines of the sub-problems. In the Poisson Boolean model of~\cite{Gou08} the analogue of the shared lines is the set of balls of radius at least~$\alpha$, which is empty with high probability; here $\cS_{kl}$ is never empty, and its law is the object the recursion has to control.

\begin{lemma}[Shared lines and aligned lines]
\label{lem:shared}
Let $R=B(y,10\alpha)$ and $R'=B(y',10\alpha)$ with $\lvert y-y'\rvert\ge70\alpha$, and $\alpha\ge1$.
\begin{enumerate}
  \item[\emph{(a)}] The lines meeting~$R$ have $\nu$-measure $20\mu\alpha$; a fortiori $\nu(\cS(R,R'))\le20\mu\alpha$.
  \item[\emph{(b)}] Let $\ell_0$ be a line meeting~$R$. The lines that pass within distance~$1$ of two points $p,q\in R$ with $\lvert p-q\rvert\ge\alpha$ of which $\ell_0$ also passes within distance~$1$ have measure at most $\beta_1\coloneqq C_1\mu/\alpha$ with $C_1\coloneqq368$.
  \item[\emph{(c)}] Let $\ell_0$ be a line meeting $R$ and~$R'$. The lines that pass within distance~$1$ of a point $p\in R$ and a point $p'\in R'$ of which $\ell_0$ also passes within distance~$1$ have measure at most $\beta_2\coloneqq C_2\mu/\alpha$ with $C_2\coloneqq1/2$.
\end{enumerate}
\end{lemma}

\begin{proof}
(a) is the perimeter formula. (b) Let $\ell$ be such a line and $\varphi$ its angle with~$\ell_0$. The points $p,q$ are within distance~$1$ of both lines and at distance at least~$\alpha$, so $\alpha\sin\varphi\le4$ and $\varphi\le(\pi/2)\sin\varphi\le2\pi/\alpha$; the normal direction of~$\ell$ lies in a set of measure at most $4\pi/\alpha$. Given the direction, $\ell$ passes within distance~$2$ of a point of the segment $\ell_0\cap B(y,11\alpha)$, whose length is at most $22\alpha$, and along this segment the distance between $\ell$ and~$\ell_0$ varies by at most $22\alpha\sin\varphi\le88$; so the offset of~$\ell$ lies in an interval of length at most $4+88=92$. The measure is at most $(\mu/\pi)(4\pi/\alpha)\cdot92=368\mu/\alpha$. (c) Now $\lvert p-p'\rvert\ge50\alpha$, so $50\alpha\sin\varphi\le4$ and $\varphi\le\pi/(25\alpha)$, the direction set has measure at most $2\pi/(25\alpha)$, the offset variation along $\ell_0\cap B(y,11\alpha)$ is at most $22\alpha\cdot4/(50\alpha)<2$, the offset interval has length at most~$6$, and the measure is at most $(\mu/\pi)(2\pi/(25\alpha))\cdot6<\mu/(2\alpha)$.
\end{proof}

The ratio $C_2/C_1=1/736$ measures how much rarer alignment across the gap between the two regions is than alignment within one region; it plays the role that the radius of a ball serving both regions plays in~\cite{Gou08}.

% =================================================================
\section{The recursion}
\label{sec:recursion}

\subsection{Abstract form}

Let $\alpha_1\ge1$. A \emph{target} is a function $f\colon[\alpha_1,\infty)\to(0,1]$, and a \emph{weight} is a measurable function $W_\alpha(x;G)\ge1$ of $\alpha\ge\alpha_1$, $x\in\R^2$ and finite sets~$G$ of lines, depending on~$G$ only through $G\cap B(x,10\alpha)$.

\begin{theorem}[Recursion]
\label{thm:recursion}
Let $(f,W)$ be a target and a weight such that
\begin{enumerate}
  \item[\emph{(B)}] \emph{(base case)} $\pi_\alpha(x;G)\le f(\alpha)W_\alpha(x;G)$ for all $\alpha\in[\alpha_1,10\alpha_1)$, $x\in\R^2$ and finite~$G$;
  \item[\emph{(R)}] \emph{(recursion)} for all $\alpha\ge\alpha_1$, $x\in\R^2$ and finite~$G$,
  \begin{equation}
  \label{eq:R}
    \sum_{k\in K}\sum_{l\in L}\E\Bigl[\min\bigl\{1,f(\alpha)W_\alpha(x_k;G\cup\cS_{kl})\bigr\}\,\min\bigl\{1,f(\alpha)W_\alpha(x_l;G\cup\cS_{kl})\bigr\}\Bigr]\le f(10\alpha)\,W_{10\alpha}(x;G)\,.
  \end{equation}
\end{enumerate}
Then $\pi_\alpha(x;G)\le f(\alpha)W_\alpha(x;G)$ for all $\alpha\ge\alpha_1$, $x\in\R^2$ and finite~$G$. In particular, if $W_\alpha(x;\varnothing)\le C_W$ for all $\alpha,x$ and $f(\alpha)\to0$, then $\Xi$ has no unbounded component almost surely.
\end{theorem}

\begin{proof}
Induction on~$n$ for $\alpha\in[10^n\alpha_1,10^{n+1}\alpha_1)$. The case $n=0$ is~(B). For $n\ge1$ write $\alpha=10\alpha'$ with $\alpha'\in[10^{n-1}\alpha_1,10^n\alpha_1)$. By~\eqref{eq:step}, the induction hypothesis applied inside the expectation with the fixed set $G\cup\cS_{kl}$, the trivial bound $\pi\le1$, and~(R),
\begin{align*}
  \pi_{10\alpha'}(x;G)&\le\sum_{k,l}\E\Bigl[\min\{1,f(\alpha')W_{\alpha'}(x_k;G\cup\cS_{kl})\}\min\{1,f(\alpha')W_{\alpha'}(x_l;G\cup\cS_{kl})\}\Bigr]\\
  &\le f(10\alpha')W_{10\alpha'}(x;G)\,.
\end{align*}
The last assertion follows from \cref{lem:perc}.
\end{proof}

The theorem separates the probabilistic input from the arithmetic: (B) is a statement about bounded regions, and (R) is an inequality between deterministic functions in which the only random object is the Poisson process $\cS_{kl}$ of shared lines, whose law is explicit. In the Poisson Boolean model the analogue of~(R) is $C\pi(\alpha)^2+\lambda C\int_\alpha^\infty r^d\nu(\dd r)\le\pi(10\alpha)$, with no weight because the dependence is carried by an event of small probability.

\subsection{The recursive weight, and the single remaining question}

Nothing prevents us from \emph{defining} $f$ and~$W$ by~(B) and~(R).

\begin{definition}[Recursive target and weight]
\label{def:recursive}
Fix $\alpha_1\ge1$. For $\alpha\in[\alpha_1,10\alpha_1)$ put $f(\alpha)\coloneqq\pi_\alpha$ and $W_\alpha(x;G)\coloneqq\max\{1,\pi_\alpha(x;G)/f(\alpha)\}$. For $\alpha\in[10^n\alpha_1,10^{n+1}\alpha_1)$, $n\ge1$, with $\alpha'=\alpha/10$, put
\begin{align}
\label{eq:frec}
  f(\alpha)&\coloneqq\sum_{k,l}\E\Bigl[\min\{1,f(\alpha')W_{\alpha'}(\alpha'k;\cS_{kl})\}\min\{1,f(\alpha')W_{\alpha'}(\alpha'l;\cS_{kl})\}\Bigr]\wedge1\,,\\
\label{eq:Wrec}
  W_\alpha(x;G)&\coloneqq\max\Bigl\{1,\ \frac1{f(\alpha)}\sum_{k,l}\E\Bigl[\min\{1,f(\alpha')W_{\alpha'}(x_k;G\cup\cS_{kl})\}\min\{1,f(\alpha')W_{\alpha'}(x_l;G\cup\cS_{kl})\}\Bigr]\Bigr\}\,,
\end{align}
where in~\eqref{eq:frec} the shared lines are those of the regions around $\alpha'k$ and $\alpha'l$ for $x=0$; by stationarity the value does not depend on~$x$.
\end{definition}

\begin{proposition}[Unconditional crossing bound]
\label{prop:unconditional}
With $f$ and $W$ of \cref{def:recursive}, $\pi_\alpha(x;G)\le f(\alpha)W_\alpha(x;G)$ for all $\alpha\ge\alpha_1$, $x$ and~$G$, and $W_\alpha(x;\varnothing)=1$ whenever $f(\alpha)<1$; in particular $\pi_\alpha\le f(\alpha)$. Consequently $\Xi$ has no unbounded component almost surely as soon as $f(\alpha)\to0$.
\end{proposition}

\begin{proof}
(B) holds by the definition at the base scales, (R) holds by~\eqref{eq:Wrec}, and for $G=\varnothing$ the maximum in~\eqref{eq:Wrec} is attained by~$1$ because the sum equals $f(\alpha)$ by~\eqref{eq:frec} when $f(\alpha)<1$. Apply \cref{thm:recursion} and \cref{lem:perc}.
\end{proof}

Everything is now in one place.

\begin{boxedclaim}{T}{The recursive target tends to zero}
\label{claim:T}
There are $\mu_0>0$ and $\lambda_0>0$ such that for $\mu\le\mu_0$, $\lambda\le\lambda_0$ and a suitable $\alpha_1$, the function~$f$ of \cref{def:recursive} satisfies $f(\alpha)\to0$ as $\alpha\to\infty$.
\end{boxedclaim}

\Cref{claim:T} is equivalent to the success of the scheme, and by \cref{prop:unconditional} it implies $\lambda_c(\mu)\ge\lambda_0$. It is a statement about an explicit sequence of expectations over Poisson processes of lines: unrolling~\eqref{eq:frec} and~\eqref{eq:Wrec}, $f(10^n\alpha_1)$ is an expectation over a tree of nested pairs of regions, the shared lines of each pair being passed to the sub-regions as fixed lines, and at the leaves of the tree the base-scale probabilities $\pi_{\alpha}(x;G)$ with $\alpha\in[\alpha_1,10\alpha_1)$ and $G$ the accumulated shared lines. We do not prove the claim. The rest of this note shows that one step of the recursion closes under a hypothesis on the weights of the shared lines (\cref{sec:onestep}), and what the hypothesis must look like to propagate (\cref{sec:discussion}).

\begin{remark}[The order of the target]
\label{rem:order}
Since $f(\alpha)\ge\pi_\alpha$, the target cannot decay faster than the crossing probability, and the following computation shows how fast that is. A crossing of the annulus from $B(0,\alpha)$ to $\partial B(0,8\alpha)$ by unit disks whose centres lie on lines requires, if all centres lie on $m$ lines running together within distance~$1$ across the annulus, a sequence of points with consecutive gaps at most~$2$ along a length of order~$\alpha$ in a Poisson process of intensity $m\lambda$, which has probability of order $\exp\{-c\alpha\,m\lambda\e^{-2m\lambda}\}$; the $m$ lines themselves cost $(\mu/\alpha)^{m}\alpha^2$ up to constants. Optimising over~$m$ gives $2m\lambda\approx\log\alpha$ and $\pi_\alpha\approx\exp\{-(\log\alpha)^2/(2\lambda)\}$ up to lower-order terms: faster than any power of~$\alpha$, and with no fixed number $j_0$ of lines. The recursion will produce a target of this type: the error term of one step is dominated by bundles of shared lines, so that $f(10\alpha)\approx10^{2-m^*(\alpha)}f(\alpha)$ with $m^*(\alpha)$ of order $\log\alpha$ (\cref{sec:onestep}), which integrates to $\exp\{-c(\log\alpha)^2\}$.
\end{remark}

\begin{remark}[Two regions suffice: a correction]
\label{rem:correction}
Our multiscale note asserted that a polynomial target forces a recursion with many well-separated sub-boxes at scales $L\mapsto L^3$, because with two sub-boxes the exponent in the basic inequality is negative. That assertion is wrong. The exponent computed there contained the positional entropy $(L/\ell)^{2a}$ of the sub-boxes, which is polynomial in~$\ell$ when $L=\ell^3$ but is the constant $\lvert K\rvert\lvert L\rvert$ at a fixed ratio; at a fixed ratio the quadratic term $f(\alpha)^2$ is below any polynomial or quasi-polynomial target $f(10\alpha)$ by a factor tending to zero, and the recursion is driven by the error term, exactly as in~\cite{Gou08}. The large number of sub-boxes was also what made the nearly complete groups of that note hard: along a row of fixed lines the sub-boxes have positional entropy $(L/\ell)^a$, which the gain of a group with few missing lines could not pay. With two regions the corresponding term is $f(\alpha)^2k_\alpha(m)^2$ against $f(10\alpha)k_{10\alpha}(m)$, and it closes for every group size (\cref{sec:discussion}).
\end{remark}

% =================================================================
\section{One recursion step under a group-weight hypothesis}
\label{sec:onestep}

In this section $\alpha\ge\alpha_1$ is fixed, $x=0$, and $f=f(\alpha)$. We show that the step from~$\alpha$ to $10\alpha$ for the target, i.e.\ the sum in~\eqref{eq:frec}, is at most $f/2$ plus an explicit \emph{across-alignment term}, provided the weights of the shared-line configurations at scale~$\alpha$ have the form one expects: groups of aligned lines are charged by their gain, with the cap~$1$.

\paragraph{Bundles and group weights.} For a region $R=B(y,10\alpha)$ and a finite set~$H$ of lines, an \emph{$m$-bundle of~$H$ in~$R$} is a set of $m$ lines of~$H$ that all pass within distance~$1$ of two points $p,q\in R$ with $\lvert p-q\rvert\ge\alpha$. Let $N_m(y;H)$ be the number of $m$-bundles of~$H$ in~$R$; $N_1(y;H)$ is the number of lines of~$H$ meeting~$R$. For $A>0$ put
\begin{equation}
\label{eq:k}
  k_\alpha(m)\coloneqq\frac{1}{\alpha^2}\Bigl(\frac{A\alpha}{\mu}\Bigr)^{m},\qquad\kappa_m\coloneqq\min\{1,f\,k_\alpha(m)\}\,.
\end{equation}
This is the gain one expects from $m$ fixed aligned lines when a crossing is carried by bundles: each fixed line replaces a random line, which by \cref{lem:shared}(b) costs a factor of order $\mu/\alpha$, so that $A$ is of the order of~$1/C_1$, and the bundle loses its freedom of position, a factor~$\alpha^2$.

\begin{definition}[Group-weight hypothesis at scale~$\alpha$]
\label{def:GW}
$\mathrm{(GW)}_\alpha$ holds with parameters $m^*\ge2$ and $A>0$ if for every $y\in\R^2$ and every finite set~$H$ of lines,
\begin{equation}
\label{eq:GW}
  \min\bigl\{1,f\,W_\alpha(y;H)\bigr\}\le f+\sum_{m=1}^{m^*}\kappa_m\,N_m(y;H)\qquad\text{with }\kappa_{m^*}\coloneqq1\,.
\end{equation}
\end{definition}

The hypothesis says that the weight of a configuration is at most the group weight, groups of $m^*$ or more lines being charged nothing beyond the cap. Put
\begin{equation}
\label{eq:pq}
  p_m\coloneqq\frac{20\mu\alpha\,\beta_1^{\,m-1}}{m!},\qquad q_{m,m',j}\coloneqq20\mu\alpha\,\beta_2^{\,j-1}\beta_1^{\,m+m'-2j}\quad(1\le j\le m\wedge m')\,.
\end{equation}

\begin{lemma}[Bundle counts]
\label{lem:counts}
Let $\cS=\cS(R_k,R_l)$, $N_m\coloneqq N_m(\alpha k;\cS)$ and $N'_{m'}\coloneqq N_{m'}(\alpha l;\cS)$. Then $\E N_m\le p_m$, and
\begin{equation}
\label{eq:pairs}
  \E[N_mN'_{m'}]\le p_mp_{m'}+\sum_{j=1}^{m\wedge m'}q_{m,m',j}\,.
\end{equation}
\end{lemma}

\begin{proof}
An $m$-bundle is an $m$-subset of~$\cS$. Ordering it, the first line meets~$R_k$, measure at most $20\mu\alpha$, and each further line passes within distance~$1$ of two points of~$R_k$ at distance at least~$\alpha$ of which the first line also passes within distance~$1$, measure at most~$\beta_1$ by \cref{lem:shared}(b); the multivariate Mecke formula~\cite[Theorem~4.4]{LP17} gives $\E N_m\le20\mu\alpha\beta_1^{m-1}/m!$. For the pairs $(T,T')$, split according to $j=\lvert T\cap T'\rvert$. The pairs with $j=0$ are pairs of disjoint subsets, and the Mecke formula for disjoint tuples gives exactly the product of the two integrals, at most $p_mp_{m'}$. For $j\ge1$, order the configuration so that a common line comes first, measure at most $20\mu\alpha$; each of the other $j-1$ common lines passes within distance~$1$ of a point of~$R_k$ and of a point of~$R_l$ of which the first line also passes within distance~$1$, because both belong to bundles in both regions, measure at most~$\beta_2$ by \cref{lem:shared}(c); and each of the $m-j$ remaining lines of~$T$ and $m'-j$ remaining lines of~$T'$ has measure at most~$\beta_1$. The number of pairs is at most the number of such ordered configurations.
\end{proof}

\begin{proposition}[One step]
\label{prop:onestep}
Assume $\mathrm{(GW)}_\alpha$ with parameters $m^*$ and~$A$. Then
\begin{equation}
\label{eq:onestep}
  \sum_{k\in K}\sum_{l\in L}\E\Bigl[\min\{1,fW_\alpha(\alpha k;\cS_{kl})\}\min\{1,fW_\alpha(\alpha l;\cS_{kl})\}\Bigr]\le\lvert K\rvert\lvert L\rvert\Bigl[\bigl(f+\Pi\bigr)^2+Q_1+\Lambda\Bigr]\,,
\end{equation}
where, with $\sigma_j\coloneqq\sum_{m=j}^{m^*}\kappa_m\beta_1^{\,m-j}$,
\begin{equation}
\label{eq:terms}
  \Pi\coloneqq\sum_{m=1}^{m^*}\kappa_mp_m,\qquad Q_1\coloneqq20\mu\alpha\,\sigma_1^2,\qquad\Lambda\coloneqq\sum_{j=2}^{m^*}20\mu\alpha\,\beta_2^{\,j-1}\sigma_j^2\,.
\end{equation}
\end{proposition}

\begin{proof}
Fix $k,l$ and use the notation of \cref{lem:counts}. By~\eqref{eq:GW} and $\E N_m\le p_m$,
\begin{align*}
  \E\Bigl[\Bigl(f+\sum_m\kappa_mN_m\Bigr)\Bigl(f+\sum_{m'}\kappa_{m'}N'_{m'}\Bigr)\Bigr]&\le f^2+2f\Pi+\sum_{m,m'}\kappa_m\kappa_{m'}\E[N_mN'_{m'}]\\
  &\le(f+\Pi)^2+\sum_{m,m'}\kappa_m\kappa_{m'}\sum_{j\ge1}q_{m,m',j}
\end{align*}
by~\eqref{eq:pairs}. For fixed~$j$, $\sum_{m,m'\ge j}\kappa_m\kappa_{m'}q_{m,m',j}=20\mu\alpha\beta_2^{j-1}\bigl(\sum_{m\ge j}\kappa_m\beta_1^{m-j}\bigr)^2=20\mu\alpha\beta_2^{j-1}\sigma_j^2$, which is $Q_1$ for $j=1$ and the $j$-th term of~$\Lambda$ for $j\ge2$. Sum over $k,l$.
\end{proof}

\begin{lemma}[Size of the terms]
\label{lem:terms}
Let $C_A\coloneqq(20/C_1)(\e^{AC_1}-1)$. Then
\begin{equation}
\label{eq:termbounds}
  \Pi\le fC_A+p_{m^*},\qquad
  Q_1\le\frac{40\,\e^{2AC_1}}{C_1^2\,\mu\alpha}\,f^2+\frac{(m^*!)^2}{10\,\mu\alpha}\,p_{m^*}^2\,.
\end{equation}
Moreover, if $\beta_1\le1/2$ and $m_0\coloneqq\min\{m\colon fk_\alpha(m)\ge1\}\wedge m^*$, then
\begin{equation*}
  \sigma_j\le m_0\,f\,\mu^{-j}\alpha^{j-2}\max_{j\le m<m_0}A^mC_1^{\,m-j}+2\beta_1^{\,m_0-j}\quad(j\le m_0),\qquad\sigma_j\le2\quad(m_0<j\le m^*)\,.
\end{equation*}
\end{lemma}

\begin{proof}
$\kappa_mp_m\le fk_\alpha(m)p_m=f\cdot20A^mC_1^{m-1}/m!$ for $m<m^*$, and $\kappa_{m^*}p_{m^*}=p_{m^*}$; summing gives the bound on~$\Pi$. For $Q_1$, $\sigma_1\le\sum_{m<m^*}fk_\alpha(m)\beta_1^{m-1}+\beta_1^{m^*-1}=f\sum_{m<m^*}A^mC_1^{m-1}/(\mu\alpha)+\beta_1^{m^*-1}\le f\e^{AC_1}/(C_1\mu\alpha)+p_{m^*}m^*!/(20\mu\alpha)$, and $(a+b)^2\le2a^2+2b^2$. For $\sigma_j$, split the sum at~$m_0$: for $m<m_0$, $\kappa_m\beta_1^{m-j}\le fk_\alpha(m)\beta_1^{m-j}=fA^mC_1^{m-j}\mu^{-j}\alpha^{j-2}$, and there are fewer than $m_0$ such terms; for $m\ge m_0$, $\kappa_m\le1$ and the geometric sum is at most $\beta_1^{(m_0-j)^+}/(1-\beta_1)\le2\beta_1^{(m_0-j)^+}$.
\end{proof}

\begin{corollary}[Halving]
\label{cor:halving}
Assume $\mathrm{(GW)}_\alpha$, $\beta_1\le1/2$, and
\begin{equation}
\label{eq:halvingconds}
  p_{m^*}\le f,\qquad f\le\frac{1}{16\lvert K\rvert\lvert L\rvert(2+C_A)^2},\qquad Q_1\le\frac{f}{8\lvert K\rvert\lvert L\rvert},\qquad\Lambda\le\frac{f}{8\lvert K\rvert\lvert L\rvert}\,.
\end{equation}
Then the sum in~\eqref{eq:onestep} is at most $f/2$, so that $f(10\alpha)\le f(\alpha)/2$ in \cref{def:recursive}.
\end{corollary}

\begin{proof}
$(f+\Pi)^2\le(f(1+C_A)+p_{m^*})^2\le f^2(2+C_A)^2\le f/(16\lvert K\rvert\lvert L\rvert)$; add the bounds on $Q_1$ and~$\Lambda$ and multiply by $\lvert K\rvert\lvert L\rvert$.
\end{proof}

\paragraph{Reading the conditions.} The first condition in~\eqref{eq:halvingconds} says that bundles large enough to be treated as free are rarer than the target; it defines the natural choice $m^*\coloneqq\min\{m\colon p_m\le f\}$, for which $p_{m^*}\le f<p_{m^*-1}$ and $m^*\approx\log(1/f)/\log(\alpha/\mu)$. The second is smallness of the target. The third holds for $\alpha\ge\alpha_0(\mu,A)$ once $p_{m^*}\le f$, because by~\eqref{eq:termbounds} $Q_1\le f\cdot[40\e^{2AC_1}f/(C_1^2\mu\alpha)+(m^*!)^2f/(10\mu\alpha)]$ and $f<p_{m^*-1}\le20\mu\alpha\beta_1^{m^*-2}=20C_1^{m^*-2}\mu^{m^*-1}\alpha^{3-m^*}$ makes the bracket at most $C^{m^*}\alpha^{2-m^*}$, small for $m^*\ge3$ and $\alpha$ large. The fourth condition is the substantial one. The dominant contributions to~$\Lambda$ come from the indices $j\ge m_0$, where the common lines are charged with $\kappa=1$: by \cref{lem:terms}, $\sum_{j\ge m_0}20\mu\alpha\beta_2^{j-1}\sigma_j^2\le160\mu\alpha\beta_2^{m_0-1}$, which is the probability that $m_0$ shared lines are aligned across both regions, up to the factor~$8$. The condition $\Lambda\le f/(8\lvert K\rvert\lvert L\rvert)$ therefore requires
\begin{equation}
\label{eq:across}
  160\,\mu\alpha\Bigl(\frac{C_2\mu}{\alpha}\Bigr)^{m_0-1}\le\frac{f}{16\lvert K\rvert\lvert L\rvert}\,,
\end{equation}
while $m_0$ is defined by $fk_\alpha(m_0)\ge1$, i.e.\ $f\ge\alpha^2(\mu/(A\alpha))^{m_0}$. Both hold together if $\alpha^2(\mu/(A\alpha))^{m_0}\ge2560\lvert K\rvert\lvert L\rvert\mu\alpha(C_2\mu/\alpha)^{m_0-1}$, that is, if $(AC_2)^{-m_0}\ge2560\lvert K\rvert\lvert L\rvert/C_2$, which is true for all $m_0\ge m_1(A)$ provided $AC_2<1$. With $A$ of the order of $1/C_1$ this is comfortably the case, and the terms $j<m_0$ of~$\Lambda$ are handled by the first bound of \cref{lem:terms} in the same way as~$Q_1$. In words: a configuration of $m_0$ shared lines aligned across both regions makes both sub-crossings free and has probability $20\mu\alpha\beta_2^{m_0-1}$; the scheme can only work if the target at scale~$\alpha$ is already at least this large, which it is once the gain constant~$A$ is below $1/C_2$, because then the same $m_0$ lines aligned within a single region, which cost $\beta_1$ per line instead of~$\beta_2$, are a configuration the target must already account for. We have not verified $\mathrm{(GW)}_\alpha$ itself; see the next section.

% =================================================================
\section{What the hypothesis must look like to propagate}
\label{sec:discussion}

\Cref{prop:onestep} concerns one step and assumes~$\mathrm{(GW)}_\alpha$ for the recursively defined weights of the configurations $\cS_{kl}$. The scheme succeeds if some hypothesis of this kind propagates: if $\mathrm{(GW)}_{\alpha}$ for all configurations implies $\mathrm{(GW)}_{10\alpha}$ for all configurations, with parameters that the recursion can tolerate. We record what the recursion~\eqref{eq:Wrec} does to the three kinds of fixed-line configurations, using $f(10\alpha)\approx10^{2-m^*}f(\alpha)$ from \cref{sec:onestep} as the order of the target at the next scale.

\paragraph{A group serving one region.} Let $G$ be an $m$-bundle in $R_k$ but not in $R_l$, $m<m^*$. The pair $(k,l)$ contributes to~\eqref{eq:Wrec} at most $f(\alpha)k_\alpha(m)\cdot f(\alpha)(1+C_A)$ up to the shared-line corrections, which is to be compared with the target's $f(10\alpha)k_{10\alpha}(m)$. The ratio is $f(\alpha)10^{m^*-2}k_\alpha(m)/k_{10\alpha}(m)\cdot(1+C_A)=f(\alpha)10^{m^*-m}(1+C_A)$, which is small since $f(\alpha)\le p_{m^*-1}\le20C_1^{m^*-2}\mu^{m^*-1}\alpha^{3-m^*}$. The group weight propagates with room to spare.

\paragraph{A group serving both regions.} Let $G$ be an $m$-bundle in both $R_k$ and~$R_l$, i.e.\ aligned across the whole scale $10\alpha$. The contribution is $f(\alpha)^2k_\alpha(m)^2$ against $f(10\alpha)k_{10\alpha}(m)$, ratio $f(\alpha)10^{m^*-2}k_\alpha(m)^2/k_{10\alpha}(m)=f(\alpha)10^{m^*-m}(A\alpha/\mu)^m\alpha^{-2}$, which with the bound on $f(\alpha)$ is at most $20C_1^{m^*-2}A^m10^{m^*-m}\mu^{m^*-1-m}\alpha^{1+m-m^*}$, small for $m\le m^*-2$ and bounded for $m=m^*-1$. This is the case that in our multiscale note required $a\ge3x/(x-2)$ sub-boxes and a separate lemma for nearly complete groups; here it closes because the positional entropy of the two regions is a constant, see \cref{rem:correction}. The case $m=m^*-1$ needs the constant~$A$ in~\eqref{eq:k} to be matched to~$C_1$, which is a matter of choosing~$A$.

\paragraph{A group making one region free.} Let $G$ be an $m^*(\alpha)$-bundle in $R_k$ of length between $\alpha$ and $10\alpha$, so that it is not an $m^*(10\alpha)$-bundle at scale $10\alpha$ and, if shorter than $10\alpha$, not a bundle at scale $10\alpha$ at all. The contribution is $1\cdot f(\alpha)(1+C_A)$, against a target $f(10\alpha)W_{10\alpha}(x;G)\approx10^{2-m^*}f(\alpha)W_{10\alpha}(x;G)$: the recursion forces $W_{10\alpha}(x;G)\ge10^{m^*-2}(1+C_A)$. The true gain of such a configuration is much smaller: a free stretch of length~$t$ in a crossing of length $70\alpha$ carried by bundles reduces the cost from $(\mu/70\alpha)^{m}$ to $(\mu/(70\alpha-t))^{m}$, a factor $\e^{mt/(70\alpha)}$, which for $t=\alpha$ and $m=m^*$ is $\e^{m^*/70}\ll10^{m^*}$. So the group weight with the bundle definition of \cref{sec:onestep} does not propagate: the two-region product bound loses the factor $f(\alpha)/f(10\alpha)\approx10^{m^*-2}$ whenever one region is free, and the weight at the larger scale has to absorb this loss although the configuration does not deserve it. Exactly the same phenomenon appeared in our multiscale note as the problem of sub-boxes made bad by tiles of fixed lines.

\paragraph{The remedy: more regions in the inclusion.} The loss comes from using only two sub-regions. If the inclusion lemma is strengthened to produce $a\ge3$ pairwise well-separated sub-regions along the crossing, which Gou\'er\'e's argument does with nets on $a$ concentric circles of radii $10,80,\dots$ scaled appropriately, then a configuration making one sub-region free still leaves $a-1\ge2$ paid sub-regions, whose product $f(\alpha)^{a-1}$ is far below $f(10\alpha)$, and the recursion does not need any weight for it. Only configurations making at least $a-1$ of the $a$ sub-regions free need to be charged, and these are bundles of $m^*$ lines covering a fraction at least $(a-1)/a$ of the crossing; their true gain is of order $(70\alpha/(\text{uncovered length}))^{m^*}\ge a^{m^*}$, which for $a\ge10$ matches or exceeds the loss $10^{m^*-2}$ of the product bound. In other words, with enough sub-regions the weight needed by the recursion for nearly complete tilings is one that the configurations actually have, and the natural candidate for the weight is the \emph{exact cost ratio}: the probability that the fixed lines, together with the cheapest random bundles filling the gaps between their aligned stretches, carry a crossing, divided by~$f$. With a fixed ratio between scales the number of sub-regions is a constant, so this costs only constants in \cref{prop:onestep}; it does not reintroduce the polynomial entropy of the multiscale note.

\paragraph{The points.} The Cox points enter only through the base case and through the caps $\min\{1,\cdot\}$: the recursion never asks how likely it is that a bundle of lines actually carries occupied disks, and a bundle of $m^*(\alpha)$ lines is treated as free. This is a loss when $m^*\lambda$ is small, since then the Cox process on the bundle has gaps larger than~$2$ at rate $m^*\lambda\e^{-2m^*\lambda}$ per unit length and the crossing of length~$7\alpha$ along the bundle is unlikely. The loss is affordable, because the recursion only needs $f\to0$ and the bundle probabilities $p_{m^*}$ already decay like $(\mu/\alpha)^{m^*}$ with $m^*\to\infty$; the parameter~$\lambda$ then enters only through the base case at scale~$\alpha_1$, where it must make $\pi_{\alpha_1}$ small. A sharper scheme would carry the point intensity on a bundle in the weight, at the price of a two-parameter weight.

\paragraph{Relation to the referee document.} The present scheme supersedes the skeleton bound, the first moment over cluster chains and the dependence step of the referee document: there is no thick cluster and no relay, since a bundle of lines released by one region is just a set of fixed lines for the next, charged by its weight. The multiscale lemmas of that document survive in simplified form, with two regions at ratio~$10$ in place of $a$ sub-boxes at ratio $L^{2/3}$, and the coverage lemma and the base case remain the tools for the two statements still open here: the base case~(B) at scale $\alpha_1$, and the propagation of a weight hypothesis, which is \cref{claim:T}.

% =================================================================
\section{Summary}
\label{sec:summary}

\begin{enumerate}
  \item Gou\'er\'e's two-scale inclusion holds for the Cox Boolean model without an error event (\cref{lem:inclusion}), with two bounded regions of radius $10\alpha$ at distance $70\alpha$ in place of the inner ball and the outer complement, and the two sub-crossings are independent given the lines meeting both regions, which become fixed lines of the sub-problems (\cref{lem:decoupling}).
  \item The recursion is a deterministic inequality~(R) for a target~$f$ and a weight~$W$, in which randomness enters only through the Poisson process of shared lines (\cref{thm:recursion}). Defining $f$ and~$W$ recursively from the base case makes the bound $\pi_\alpha\le f(\alpha)$ unconditional (\cref{prop:unconditional}), and the subcritical phase follows from the single statement $f(\alpha)\to0$ (\cref{claim:T}).
  \item One step of the recursion closes, $f(10\alpha)\le f(\alpha)/2$, if the weights of shared-line configurations at scale~$\alpha$ are bounded by group weights with bundles of size $m^*(\alpha)$ treated as free (\cref{prop:onestep}); the configurations aligned across both regions are controlled by the factor $C_2/C_1$ between alignment across the gap and alignment within a region.
  \item Group weights propagate for groups serving one or both regions, which removes the nearly complete groups and the many sub-boxes of the multiscale note (\cref{rem:correction}), but not for groups making one region free. The remedy is an inclusion with a constant number $a\ge3$ of sub-regions and an exact-cost weight for nearly complete tilings; carrying this out, together with the base case, is what remains.
\end{enumerate}

\section*{Acknowledgements}

BJ acknowledges the financial support of the Leibniz Association within the Leibniz Junior Research Group on Probabilistic Methods for Dynamic Communication Networks (DYCOMNET) as part of the Leibniz Competition (grant no.~J105/2020), of the Berlin Mathematics Research Center MATH+ through the project EF-MA-Sys-2 on Information Flow \& Emergent Behavior in Complex Networks, and of the Deutsche Forschungsgemeinschaft (DFG) through Project P27 within the SPP~2265.

\section*{Use of AI tools}

The proofs, computations and the text of this note were developed with substantial assistance of the AI systems Claude (Anthropic) and ChatGPT Astra~6 (OpenAI), in an interactive process directed by the authors.

\bibliographystyle{amsalpha}
\bibliography{references}

\end{document}
